\href{https://onlinelibrary.wiley.com/doi/book/10.1002/9780470061602}{Encyclopedia of Quantitative Finance}
\href{https://onlinelibrary.wiley.com/doi/book/10.1002/9781118182635}{Encyclopedia of Financial Models}


\subsection{Portfolio Allocation and Short Selling}

\noindent We will use the following notation. An element \(\xbar{x}\) of \(\R^{d+1}\) is a vector
\begin{equation} \label{eq:VectorDefinition}
    \xbar{x} = \left(x^{(0)}, x^{(1)}, \dots, x^{(d)}\right)
\end{equation}
made of \(d+1\) components. The scalar product \(\xbar{x} \cdot \xbar{y}\) of two vectors \(\xbar{x}, \xbar{y} \in \R^{d+1}\) is defined by
\begin{equation} \label{eq:ScalarProductDefinition}
    \langle\xbar{x}, \xbar{y}\rangle \coloneqq x^{(0)}y^{(0)} + x^{(1)}y^{(1)} + \dots + x^{(d)}y^{(d)}.
\end{equation}
The vector
\begin{equation} \label{eq:PriceVectorT0}
    \xbar{S}_0 = \left(S_0^{(0)}, S_0^{(1)}, \dots, S_0^{(d)}\right)
\end{equation}
denotes the prices at time \(t=0\) of \(d+1\) assets. Namely, \(S_0^{(i)} > 0\) is the price at time \(t=0\) of asset \(n^o\) \(i=0, 1, \dots, d\).

The asset values \(S_1^{(i)} > 0\) of assets No \(i=0, 1, \dots, d\) at time \(t=1\) are represented by the vector
\begin{equation} \label{eq:PriceVectorT1}
    \xbar{S}_1 = \left(S_1^{(0)}, S_1^{(1)}, \dots, S_1^{(d)}\right),
\end{equation}
whose components \(\left(S_1^{(1)}, \dots, S_1^{(d)}\right)\) are random variables defined on a probability space \((\Omega, \cF, \Pm)\).

In addition, we will assume that asset \(n^o\) 0 is a risk-less asset (of savings account type) that yields an interest rate \(r > 0\), \textit{i.e.} we have
\begin{equation} \label{eq:RisklessAssetDynamics}
    S_1^{(0)} = (1+r)S_0^{(0)}.
\end{equation}

\begin{defn}{Portfolio}{PortfolioDef}
    A portfolio \textit{based on the assets \(0,1,\dots,d\)} is a vector
    \begin{equation} \label{eq:PortfolioVector}
        \xbar{\xi} = \left(\xi^{(0)}, \xi^{(1)}, \dots, \xi^{(d)}\right) \in \R^{d+1},
    \end{equation}
    in which \(\xi^{(i)}\) represents the (possibly fractional) quantity of asset \(n^o\) \(i\) owned by an investor, \(i=0, 1, \dots, d\).
\end{defn}

\noindent The \textit{price} of such a portfolio, or the cost of the corresponding investment, is given by
\begin{equation} \label{eq:PortfolioPriceT0}
    \xbar{\xi} \cdot \xbar{S}_0 = \sum_{i=0}^d \xi^{(i)} S_0^{(i)} = \xi^{(0)} S_0^{(0)} + \xi^{(1)} S_0^{(1)} + \dots + \xi^{(d)} S_0^{(d)}
\end{equation}
at time \(t=0\). At time \(t=1\), the \textit{value} of the portfolio has evolved into
\begin{equation} \label{eq:PortfolioValueT1}
    \xbar{\xi} \cdot \xbar{S}_1 = \sum_{i=0}^d \xi^{(i)} S_1^{(i)} = \xi^{(0)} S_1^{(0)} + \xi^{(1)} S_1^{(1)} + \dots + \xi^{(d)} S_1^{(d)}.
\end{equation}

\noindent There are various ways to construct a portfolio allocation \((\xi^{(i)})_{i=0,1,\dots,d}\).

\begin{enumerate}[label=\roman*), leftmargin=*, parsep=6pt]
    \item If \(\xi^{(0)} > 0\), the investor puts the amount \(\xi^{(0)} S_0^{(0)} > 0\) on a savings account with interest rate \(r\).

    \item If \(\xi^{(0)} < 0\), the investor borrows the amount \(-\xi^{(0)} S_0^{(0)} > 0\) with the same interest rate \(r\).

    \item For \(i = 1, 2, \dots, d\), if \(\xi^{(i)} > 0\), the investor purchases a (possibly fractional) quantity \(\xi^{(i)} > 0\) of the asset \(n^o\) \(i\).

    \item If \(\xi^{(i)} < 0\), the investor borrows a quantity \(-\xi^{(i)} > 0\) of asset \(i\) and sells it to obtain the amount \(-\xi^{(i)} S_0^{(i)} > 0\).
\end{enumerate}

\noindent In the latter case one says that the investor \textit{short sells} a quantity \(-\xi^{(i)} > 0\) of the asset \(n^o\) \(i\), which lowers the cost of the portfolio.

\begin{defn}{Short Selling Ratio}{ShortSellingRatioDef}
    The \textit{short selling ratio}, or percentage of daily turnover activity related to short selling, is defined as the ratio of the number of daily short sold shares divided by daily volume.
\end{defn}

\noindent Profits are typically generated by first buying at a low price and then selling at a higher price. Short sellers apply the same rule, but in reverse chronological order: first sell high and then buy low, if possible, by following the procedure below.

\begin{enumerate}
    \item Borrow the asset \(n^o\) \(i\).

    \item At time \(t=0\), sell the asset \(n^o\) \(i\) on the market at the price \(S_0^{(i)}\) and invest the amount \(S_0^{(i)}\) at the interest rate \(r > 0\).

    \item Buy back the asset \(n^o\) \(i\) at time \(t=1\) at the price \(S_1^{(i)}\), with hopefully \(S_1^{(i)} < (1+r)S_0^{(i)}\).

    \item Return the asset to its owner, with possibly a (small) fee \(p > 0\).\footnote{The cost \(p\) of short selling will not be taken into account in later calculations.}
\end{enumerate}

\noindent At the end of the operation, the profit made on share \(n^o\) \(i\) equals
\begin{equation} \label{eq:ShortSellingProfit}
    (1+r)S_0^{(i)} - S_1^{(i)} - p > 0,
\end{equation}
which is positive provided that \(S_1^{(i)} < (1+r)S_0^{(i)}\) and \(p > 0\) is sufficiently small.

\subsection{Arbitrage}


\subsection{Risk-neutral Probability Measures}


\subsection{Hedging of Contingent Claims}
\begin{defn}{Contingent Claim}{ContingentClaim}
A contingent claim is a financial derivative whose payoff \(C : \Omega \longrightarrow \mathbb{R}\) is a random variable depending on the realisation of uncertain events.
\end{defn}

\subsection{Market Completeness}